\documentclass[11pt]{article}
\usepackage[utf8]{inputenc}
\usepackage{algorithm}
\usepackage{amsmath}
\usepackage{graphicx}
\usepackage{caption}
\usepackage{subcaption}
\usepackage{mathtools}
\usepackage{amsfonts}
\usepackage{hyperref}
\usepackage{geometry}
\usepackage{multicol}
\usepackage{capt-of}
\usepackage{url}
\usepackage{chemmacros}
\usepackage{upgreek}
\usepackage{isomath}
\usepackage{cancel}
\usepackage{bm}
\usepackage[italian]{babel}
\geometry{a4paper, top=2.5cm, bottom=2.5cm, left=3cm, right=3cm, bindingoffset=5mm}
\usepackage[suftesi]{frontespizio}
\usepackage{siunitx}
\usepackage{float}
\usepackage{verbatim}

\title{Tutta la statistica da sapere per Fenomenologia del Modello Standard delle Particelle Elementari\\ }
\author{Nicolò Favagrossa\\nicolo.favagrossa@studenti.unimi.it}
\date{\today}
\begin{document}
\maketitle
\newpage
\tableofcontents
\newpage
\section{Introduzione}
Il più semplice esempio di intervalli di confidenza sono le "barre di errore". Al primo anno del laboratorio di fisica abbiamo imparato a mettere le misure in questa maniera
\begin{equation}
    \text{misura} = \text{valore} \pm \text{errore}
\end{equation}
che significa che è la miglior stima della $X$ cin la sua incertezza $\sigma_X$ per una certa quantità. Il messaggio è: "sei confidente che il valore vero $X^\text{true}$ giaccia nell'intervallo $[X-\sigma_X, X+\sigma_X]$ con una probabilità ragionevole".
Più precisamente stiamo assumento una fluttuazione Gaussiana dei nostri dati e crediamo che:
\begin{equation}
    X^\text{true} \in [X-\sigma_X, X+\sigma_X] \quad \text{con una probabilità del } 68\%
\end{equation}
questa è chiamata livello di confidenza e $\left[ X-\sigma_X\;X+\sigma_X \right]$ è chiamato intervallo di confidenza al $68\%$.
Attenzione però! $X^\text{True}$ non è conosciuto! La probabilità quindi si riferisce a $X$ e $\sigma_X$ che fluttuano da misura a misura.

Ora però dobbiamo generalizzare queso concetto in questi seguenti casi:\
\begin{itemize}
    \item le fluttuazioni non sono necessariamente Gaussiane
    \item il livello di confidenza voluto potrebbe essere diverso dal $68\%$ (ad esempio $95\%$)
    \item potremmo essere interessati a più di una quantità (ad esempio $X$ e $Y$) che potrebbero essere anche correlate
    \item la stessa quantità potrebbe essere misurata da più di un esperimento con diversi metodi, data sets, etc...
\end{itemize}
\begin{figure}[h!]
    \centering
    \includegraphics[width=0.9\textwidth]{immagini/Confidence_Intervals.png}
    \label{fig:ff}
\end{figure}
Cosa significano questi plot? Questi sono intervalli di confidenza costruiti dai dati $\Rightarrow$ stocastici. Hanno una probabilita uguale al Confidence Level (CL) di contenere ("ricoprire") il valore vero (sconosciuto). Come li costruisco? lo vedremo in seguito!

Un primo esempio è sicuramente l'accoppiamento del bosone di Higgs ai fermioni e ai bosoni di Gauge:
\begin{equation}
    g_f^{SM}=\frac{m_f}{v} \quad \text{e} \quad g_V^{SM}=2\frac{m_V^2}{v}
\end{equation}
E noi misuriamo $g_f$ e $g_V$ e vogliamo sapere se sono compatibili con il modello standard. Quindi vogliamo costruire un intervallo di confidenza per $g_f$ e $g_V$ e vedere se il punto $g_f^{SM}$, $g_V^{SM}$ è incluso in questo intervallo.


\end{document}

